% =============================================================================
%   III. DATA, INDEX CONSTRUCTION AND VALIDATION DESIGN
% =============================================================================
\section{Data and Index Construction}
\label{sec:data}

\subsection{Universe and Snapshot}
\label{sec:universe}

The candidate universe contains \NRawCompanies firms: mid-cap constituents
of the S\&P MidCap~400, the Russell Midcap and the NIFTY Midcap~150
\cite{nse_midcap150}, plus \NLarge S\&P~500 large caps that serve only as a
reference group in normalisation. Fundamentals and trailing market statistics
come from Yahoo Finance \cite{yahoo_finance, yfinance} and SEC EDGAR
\cite{sec_edgar} and are dated at the close of \SnapshotDate, which we
verified by matching every cached last price to the exchange close. A firm
must have traded on the snapshot date to be scored. This filter removes
\NDroppedNotTrading firms (\DroppedNotTrading) that had already been acquired
or delisted and whose inclusion would have let the evaluation window
``see'' corporate events that were already known. The scored universe has
\NScored firms, of which \NMid are mid-caps (\NUS US, \NIndia India) in
\NSectors sectors; all analysis concerns the mid-caps. Indian monetary fields
are converted at the \SnapshotDate close of INR/USD (\FXRate).

\subsection{Public-Data ESG with Cell-Level Provenance}
\label{sec:esg}

The ESG schema declares \NIndDeclared indicators (\NIndDeclaredE E,
\NIndDeclaredS S, \NIndDeclaredG G). Each company--indicator cell is filled
by the first available source in a fixed order, and the source is recorded.
The order runs from direct measurement to imputation: the EPA Toxics Release
Inventory (TRI) for US facilities; ISS governance QualityScores redistributed
by Yahoo Finance (audit, board, compensation, shareholder-rights and overall
risk on a 1--10 scale, stored as $(10-r)\times 10$ so that higher is
better); SEC XBRL filing items; Yahoo ownership and analyst-coverage fields;
membership of public initiatives (CDP, SBTi, RE100, UN Global Compact);
within-sector percentiles of financial ratios, each with a published
rationale; and finally the sector or global median. No random noise is added
at any stage. Appendix~\ref{app:indicators} gives the source of each
indicator.

Table~\ref{tab:provenance} shows the outcome. Only \NIndUsed indicators vary
across firms. No mid-cap matched an EPA facility record, and the
Sustainalytics pillar scores that Yahoo once exposed are no longer
populated, so the E and S pillars consist almost entirely of financial
ratios and proxies. The G pillar draws mainly on ISS scores (\ISSShareG\% of
cells, available for \NISSFirms of \NMid firms), followed by ownership and
analyst-coverage fields.

\begin{table}[t]
\centering
\caption{Provenance of ESG cells used in scoring (\% of cells per pillar,
mid-cap universe). ``Used/declared'' counts indicators with cross-sectional
variation.}
\label{tab:provenance}
\small
% Auto-generated by scripts/26_paper_artifacts.py
\begin{tabular}{@{}lrrr@{}}
\toprule
Source of cell & E & S & G \\
\midrule
Indicators used / declared & 4/11 & 7/11 & 10/12 \\
\midrule
ISS governance score & -- & -- & 39 \\
Yahoo ownership/coverage field & -- & -- & 28 \\
SEC XBRL item & -- & -- & 3 \\
Public registry (0/1) & 25 & 14 & -- \\
Financial ratio & -- & 14 & -- \\
Financial proxy & 75 & 70 & 12 \\
Imputed & \ensuremath{<}1 & 1 & 19 \\
\bottomrule
\end{tabular}

\end{table}

\subsection{Factor Scores}
\label{sec:factors}

Continuous variables are winsorised at the 1st/99th percentile (2.5th/97.5th
for valuation ratios) and standardised with a robust $z$-score based on the
median absolute deviation (MAD),
$z=(x-\operatorname{med}(x))/(1.4826\,\mathrm{MAD}(x))$, clipped to $[-3,3]$
\cite{iglewicz_1993}. For the operational, risk-adjusted, value, growth and
stability factors, indicators are first converted to within-sector
percentiles, and indicators that fail a Shapiro--Wilk normality test
($p<0.05$) receive normal scores. The financial, operational, growth, value
and stability scores are then blended 70:30 with their own within-sector
percentiles. Lower-is-better indicators are sign-flipped only when they are
stored as raw magnitudes; Section~\ref{sec:audit} shows why this matters.

Within the ESG block, indicators are averaged within categories (e.g.\
emissions, labour, board), categories are weighted into E, S and G, and the
pillars are combined with sector-specific SASB materiality weights
(Appendix~\ref{app:sasb}) \cite{khan_2016, sasb_2023}. Governance values are
quantile-harmonised by source before the G pillar is formed, so that ISS and
SEC-derived inputs share a scale.

The eight financial factors are conventional. The financial factor covers
profitability, margins, cash and payout; the operational factor revenue per
employee, R\&D intensity, free cash flow (FCF), gross margin and headcount;
the risk-adjusted factor one-year Sharpe and Sortino ratios, maximum
drawdown, Amihud illiquidity \cite{amihud_2002}, skewness and kurtosis; the
growth factor revenue, earnings and sustainable growth; the value factor
inverted P/E, P/B, P/S, EV/EBITDA and EV/revenue; and the stability factor
liquidity, leverage and 30-day volatility. A sector-position factor measures
size and efficiency relative to sector peers. A market factor (liquidity,
beta, volatility and 1- to 6-month momentum) is computed but receives zero
weight in every profile, so that firms are not scored on the trailing
returns that earlier versions of the pipeline used as the validation target.

\subsection{Aggregation and Investor Profiles}
\label{sec:aggregation}

For profile $p$ with weights $w^{(p)}$, the preference score is a weighted
mean of factor percentiles,
\begin{equation}
C^{(p)}_i=\sum_k w^{(p)}_k\,\frac{100\,\operatorname{rank}(F_{ik})}{N},
\qquad \sum_k w^{(p)}_k=1,
\label{eq:composite}
\end{equation}
with average ranks for ties and median fill for missing factors. Ranks are
taken over the $N=\NScored$ scored firms, so the \NLarge large caps act as a
fixed reference group. Table~\ref{tab:weights} lists the weights of the three
profiles; the balanced profile puts \WBalESG\% on ESG. The index portfolio is
the equal-weighted set of the top
$\lceil \TopPct\%\cdot N_{\mathrm{mid}}\rceil=\TopN$ mid-caps
($N_{\mathrm{mid}}=\NMid$).

The profile weights were set by the authors from in-sample rank correlations
with \emph{trailing} momentum, a practice we criticise in
Section~\ref{sec:res_insample}, and were last revised in July~2026, after the
snapshot. The OOS analysis therefore also reports the February~2026 weights,
which predate the snapshot, and equal weights.

\begin{table}[t]
\centering
\caption{Effective profile weights (\%) after removing the market factor.}
\label{tab:weights}
\small
% Auto-generated by scripts/26_paper_artifacts.py
\begin{tabular}{@{}lrrr@{}}
\toprule
Factor & ESG-first & Balanced & Financial-first \\
\midrule
ESG composite & 56.1 & 23.2 & 5.6 \\
Financial & 8.2 & 18.9 & 27.8 \\
Risk-adjusted & 8.2 & 15.8 & 20.0 \\
Growth & 6.1 & 14.7 & 16.7 \\
Stability & 6.1 & 8.4 & 7.8 \\
Sector position & 7.1 & 8.4 & 5.6 \\
Operational & 5.1 & 6.3 & 8.9 \\
Value & 3.1 & 4.2 & 7.8 \\
Market & 0.0 & 0.0 & 0.0 \\
\bottomrule
\end{tabular}

\end{table}

\subsection{Audit Corrections}
\label{sec:audit}

Before any OOS number was computed, we audited the pipeline and corrected
its errors (Appendix~\ref{app:audit} gives the full list). The most
consequential was a double direction flip: the Scope~1 emissions and ESG
risk-rating indicators were already stored as ``higher is better'' scores
and were flipped again, which rewarded high-emission and high-governance-risk
firms. The Spearman correlation between the pre- and post-audit E pillars is
\RhoPrePostE, and for G it is \RhoPrePostG. The audit also removed
multiplicative Gaussian noise from median imputation, a ``monotonicity
correction'' that fitted each factor by isotonic regression to the proxy
later used to validate it, an undocumented gate that set the bottom 30\% of
firms by financial score to zero, an exponent of 1.5 on preference scores,
and a live, non-snapshot FX rate. After these corrections the balanced
composite has a rank correlation of \RhoPrePostBal with the pre-audit
version, but only \OverlapPrePost of the \TopN top names are shared.

% -----------------------------------------------------------------------------
\section{Validation Design}
\label{sec:validation}

\subsection{In-Sample Diagnostics}
\label{sec:insample_design}

On the snapshot cross-section we report conventional diagnostics:
variance-inflation factors (VIFs), principal component analysis (PCA) and
agreement between profiles. We also re-run the pipeline's former validation
designs to measure how much they inflate apparent quality. These are rank
information coefficients (ICs) against 1-, 3- and 6-month \emph{trailing}
returns, which end on the snapshot date; ICs against a ``forward quality
proxy'', which is in fact a \emph{contemporaneous} composite of ROA, ROE,
margins, growth and liquidity that shares inputs with several factors; and
quintile monotonicity tests against the same proxy.

\subsection{Out-of-Sample Test}
\label{sec:oos_design}

Scores are fixed at the \SnapshotDate close, and outcomes are measured over
(\SnapshotDate,\,\OOSEndDate], which spans \TradingDaysUS US and
\TradingDaysIndia Indian trading days. From split- and dividend-adjusted
closes we compute returns at one and three months and over the full window,
together with realised volatility, maximum drawdown and beta against the
S\&P~500 or NIFTY~50. Returns are demeaned within country, because the two
markets had different index returns (\IdxRetUS\% and \IdxRetIndia\%);
sector-and-country neutralisation is a robustness check. \NExited firms were
acquired during the window (\ExitedTickers) and have no end price, so we
report bounds that assign them the 5th or 95th percentile return.

Three hypotheses were fixed before outcomes were computed. All are
one-sided, with Holm correction across the family \cite{holm_1979}.
\begin{itemize}
\item[\textbf{H1}] The balanced composite has a positive rank IC with the
  country-neutral full-window return.
\item[\textbf{H2}] Given trailing (2-year) volatility, higher ESG predicts lower
  realised OOS volatility (partial Spearman).
\item[\textbf{H3}] The top-\TopN balanced portfolio beats the equal-weighted
  universe.
\end{itemize}
Confidence intervals for H1--H3 come from 5{,}000 firm-level bootstrap
resamples (1{,}000--2{,}000 for secondary analyses), and in H3 the portfolio
is re-formed in each resample. Secondary analyses cover factor ICs by
horizon (Holm-corrected within horizon), portfolios by profile, and
ESG--risk partial correlations and rank regressions with sector and country
fixed effects and heteroskedasticity-robust (HC3) errors.

With $N=\NOOS$, a two-sided test at $\alpha=0.05$ has 80\% power only for
$|\rho|\ge\MDEOOS$. ICs of the size usually reported for single factors
($0.02$--$0.10$ \cite{grinold_kahn_2000}) are therefore not detectable in one
cross-section, and a null for H1 means ``no large effect'' rather than ``no
effect''.

The window had already ended when the hypotheses were written down, and the
design was not externally pre-registered \cite{nosek_2018}. Its safeguards
are procedural: universe, weights and scores are frozen artefacts dated at
the snapshot, post-snapshot prices were downloaded only after the hypotheses
and code were fixed, and all alternative weightings are reported.

\subsection{Secondary Analyses: Is the Risk Signal ESG?}
\label{sec:ext_design}

If H2 holds, three alternative explanations must be ruled out before the
effect can be attributed to ESG. These tests are exploratory and were
specified after H1--H3 had been evaluated.

The first is that ESG proxies for a known low-risk characteristic
\cite{ang_2006, frazzini_pedersen_2014, asness_2019}. We regress the
percentile rank of realised volatility on the ESG rank and on the ranks of
trailing volatility, trailing beta, size, quality (profitability score),
leverage (debt/equity), liquidity (log dollar volume), six-month momentum and
value (price/book), with sector and country fixed effects and HC3 errors
(\emph{H4}). We also ask whether ESG improves an out-of-fold volatility
forecast (\CVFolds-fold cross-validation, \CVRepeats repeats), and we
double-sort firms on trailing volatility (terciles within country) and ESG.

The second is that the signal comes from the financial proxies rather than
from measured ESG data. For each firm we average the percentile ranks of its
indicator cells separately by source: \emph{measured} (ISS, SEC, public
registry; \CellsMeas cells per firm on average), \emph{proxy}
(financial-statement proxies; \CellsProxy) and \emph{imputed} (sector or
global medians; \CellsImp). We then re-run the risk tests on each sub-score.

The third is that ESG is a slow-moving record of past volatility. We split
the window at \SplitDate and ask whether ESG predicts second-half volatility
given \emph{first-half} realised volatility, which the scores never saw but
an investor would have.

Firms whose realised volatility fell below a quarter of their trailing
volatility (\PinnedTickers) were trading at announced takeover prices. We
report results without them, by country, with winsorised volatility, and
without utilities and real estate. For the return null, we report the OOS
IC over \WDraws weight vectors drawn uniformly from the simplex and over a
sweep of the ESG weight from 0 to 1, month-by-month ICs, and top-\TopN
excess returns net of round-trip costs of \CostBpsUS~bps (US) and
\CostBpsIndia~bps (India, including securities transaction tax), and of
three times those costs.

\subsection{Pre-Registered Second Window}
\label{sec:prereg_design}

To make the next test confirmatory, we froze the \NMid firm scores and a
four-hypothesis analysis plan on \PreregFrozen, before the second window
(\PreregAnchor,\,\PreregEnd] opened. The plan, the frozen scores and the
evaluation code are hashed (SHA-256), and the package was pushed to the
public project repository on \PreregFrozen; that push is its timestamp. At
the time of writing the window is still open and no outcome in it has been
computed. The plan keeps H1--H3, adds H4 (Section~\ref{sec:ext_design}) and
applies the Holm correction across all four. In H2 and H4 the controls are
those observable at the anchor date, namely six months of realised
volatility, beta and return. The evaluation code refuses to run before the
window closes or if any registered file has changed. A dry run of the same
code on (\DryAnchor,\,\DryEnd] is included in the registration.

\subsection{Rank Uncertainty}
\label{sec:rank_design}

Rank intervals combine two sources of uncertainty over $B=\BootN$ draws. The
reference universe used for percentiles is resampled with replacement, and
the profile weights are drawn from a Dirichlet($\kappa w$) with
$\kappa=\BootKappa$, about $\pm20\%$ relative variation for a weight of 0.2.
The re-implementation reproduces the deployed composite with Spearman
$\rho=\BootRecon$.
